\chapter{Metadata strony}

\section{Po co w ogole sa metadata}
Kiedy otwierasz strone w przegladarce, widzisz glownie to, co trafia do sekcji
widocznej dla uzytkownika. Ale dokument HTML ma tez \mc{<head>}, a w nim informacje,
ktore nie zawsze sa od razu widoczne, ale maja ogromne znaczenie praktyczne.
Sa to m.in.:

\begin{itemize}
  \item tytul strony na karcie przegladarki,
  \item opis dla wyszukiwarek,
  \item wskazowki dla robotow indeksujacych,
  \item dane dla social media,
  \item linki do arkuszy stylow, ikon i manifestu,
  \item skrypty dolaczane do dokumentu.
\end{itemize}

MoonChunk pozwala definiowac te informacje bezposrednio w kodzie \mc{.mncnk}, bez
koniecznosci edytowania zewnetrznego szablonu HTML.

\section{Podstawowa skladnia metadata}
Metadata zapisuje sie jako proste instrukcje typu \mc{nazwa: wartosc;}. Przyklad:

\begin{lstlisting}[style=moonchunk,caption={Najprostsze metadata strony}]
chunk "Main" {
  output: "./dist";
  title: "MoonChunk Docs";
  metaDescription: "Dokumentacja uzytkownika MoonChunk.";

  page "" {
    content {
      <h1>MoonChunk</h1>
    };
  };
};

moon(Main);
\end{lstlisting}

W tym przykladzie \mc{title} ustawia tytul dokumentu, a \mc{metaDescription} opis.

\section{Dlaczego to rozwiazanie jest wygodne}
W klasycznym podejsciu trzeba byloby otwierac szablon HTML i recznie wstawiac wiele
znacznikow do \mc{<head>}. MoonChunk robi to za Ciebie. Ty pracujesz na poziomie
intencji: \emph{jaki ma byc tytul}, \emph{jaki ma byc opis}, \emph{jaka ma byc ikona},
a jezyk wstawia to we wlasciwe miejsca wewnetrznego szablonu dokumentu.

To sprawia, ze:

\begin{itemize}
  \item pliki \mc{.mncnk} staja sie glownym miejscem konfiguracji strony,
  \item latwiej utrzymac spojnosc metadata,
  \item prosciej jest nadpisac tylko pojedyncze pola na konkretnej stronie.
\end{itemize}

\section{Metadata wspolne i lokalne}
Bardzo czesty scenariusz wyglada tak:

\begin{itemize}
  \item kilka pol metadata jest wspolnych dla calego serwisu,
  \item wybrane pola, np. \mc{title} albo \mc{metaDescription}, roznia sie per strona.
\end{itemize}

MoonChunk obsluguje taki model bardzo naturalnie. Mozesz utworzyc osobny chunk z
metadata wspolnymi, a potem dolaczyc go przez \mc{@include}.

\begin{lstlisting}[style=moonchunk,caption={Wspolny chunk z metadata}]
export chunk "Metadata" {
  lang: "pl";
  charset: "utf-8";
  viewport: "width=device-width, initial-scale=1";
  metaAuthor: "MoonChunk Team";
  metaRobots: "index,follow";
  themeColor: "#0f172a";
};
\end{lstlisting}

A nastepnie:

\begin{lstlisting}[style=moonchunk,caption={Nadpisywanie wybranych metadata w konkretnym chunku}]
chunk "DocsHome" {
  @include Metadata;
  output: "./dist";
  title: "Dokumentacja MoonChunk";
  metaDescription: "Strona glowna dokumentacji.";

  page "" {
    content {
      <h1>Dokumentacja</h1>
    };
  };
};
\end{lstlisting}

Tutaj wspolne pola pochodza z \mc{Metadata}, a pola lokalne doprecyzowuja dana strone.
W praktyce ostatnia ustawiona wartosc metadata wygrywa.

\section{Kategorie wspieranych metadata}
MoonChunk wspiera szeroki zestaw pol. Najwygodniej myslec o nich w grupach.

\subsection*{Metadata dokumentu i jezyka}
To pola opisujace podstawowy kontekst strony:

\begin{itemize}
  \item \mc{lang},
  \item \mc{dir},
  \item \mc{htmlClass},
  \item \mc{charset},
  \item \mc{viewport}.
\end{itemize}

\subsection*{SEO i informacje podstawowe}
To pola najczesciej wykorzystywane przez wyszukiwarki:

\begin{itemize}
  \item \mc{title},
  \item \mc{metaDescription},
  \item \mc{metaKeywords},
  \item \mc{metaAuthor},
  \item \mc{metaRobots},
  \item \mc{themeColor}.
\end{itemize}

\subsection*{Linki zasobow}
To pola odpowiadajace za zewnetrzne zasoby dokumentu:

\begin{itemize}
  \item \mc{canonicalUrl},
  \item \mc{faviconHref},
  \item \mc{appleTouchIconHref},
  \item \mc{manifestHref},
  \item \mc{preloadLinks},
  \item \mc{preconnectLinks},
  \item \mc{styles},
  \item \mc{headScripts}.
\end{itemize}

\subsection*{Open Graph i Twitter}
Te pola sa przydatne przy udostepnianiu stron w social media:

\begin{itemize}
  \item \mc{ogType}, \mc{ogTitle}, \mc{ogDescription}, \mc{ogImage}, \mc{ogUrl},
        \mc{ogSiteName}, \mc{ogLocale},
  \item \mc{twitterCard}, \mc{twitterSite}, \mc{twitterCreator}, \mc{twitterTitle},
        \mc{twitterDescription}, \mc{twitterImage}.
\end{itemize}

\subsection*{Elementy struktury calego dokumentu}
MoonChunk obsluguje takze pola, ktore trafiaja do ciala strony poza samym \mc{content}:

\begin{itemize}
  \item \mc{bodyClass},
  \item \mc{pageId},
  \item \mc{topBar},
  \item \mc{header},
  \item \mc{footer},
  \item \mc{modals},
  \item \mc{scripts},
  \item \mc{bodyEndExtra}.
\end{itemize}

\section{Co sie dzieje, gdy pole jest puste}
MoonChunk ma wewnetrzny szablon bazowy dokumentu. Gdy niektore opcjonalne pola metadata
pozostana puste, odpowiadajace im puste znaczniki nie musza zostac zachowane w finalnym
HTML-u. To wazne, bo dzieki temu wynikowy dokument nie jest zasmiecany pustymi
\mc{meta}, \mc{link} albo pustym \mc{title}.

W praktyce warto jednak myslec tak: jesli dana informacja ma znaczenie dla strony,
uzupelnij ja swiadomie. Jesli nie ma znaczenia, mozesz pozostawic pole puste.

\section{Metadata a tresc strony}
Czesta pomylka poczatkujacych polega na mieszaniu tych dwoch warstw. Na przyklad:

\begin{itemize}
  \item \mc{title} nie zastępuje widocznego naglowka \mc{<h1>},
  \item \mc{metaDescription} nie zastępuje pierwszego akapitu na stronie,
  \item \mc{styles} nie zastępuje samego pliku CSS.
\end{itemize}

Innymi slowy:

\begin{itemize}
  \item metadata opisują dokument i jego otoczenie,
  \item \mc{content} opisuje glowna tresc widoczna dla uzytkownika.
\end{itemize}

Dobra strona zwykle potrzebuje jednego i drugiego.

\section{Przyklad rozsadnego zestawu startowego}
Dla wielu projektow sensowny zestaw bazowy moze wygladac tak:

\begin{lstlisting}[style=moonchunk,caption={Rozsadny zestaw metadata startowych}]
chunk "Main" {
  output: "./dist";
  lang: "pl";
  charset: "utf-8";
  viewport: "width=device-width, initial-scale=1";
  title: "MoonChunk";
  metaDescription: "Przyklad strony wygenerowanej w MoonChunk.";
  metaAuthor: "Twoje Imie";
  metaRobots: "index,follow";
  themeColor: "#111827";
  styles: "<link rel='stylesheet' href='/assets/app.css'>";

  page "" {
    content {
      <h1>MoonChunk</h1>
      <p>Start projektu.</p>
    };
  };
};
\end{lstlisting}

To juz daje porzadny punkt wyjscia do rozwijania serwisu.

\section{Kiedy wydzielic metadata do osobnego chunku}
Jesli kopiujesz te same pola do kilku miejsc, to prawie zawsze sygnal, ze warto
wydzielic wspolny \mc{chunk} z metadata. Najczesciej robi sie to wtedy, gdy:

\begin{itemize}
  \item masz wiecej niz jedna strone,
  \item kazda strona ma ten sam \mc{lang}, \mc{charset}, \mc{viewport},
  \item wspolne sa ikony, skrypty i style,
  \item tylko tytul i opis zmieniaja sie lokalnie.
\end{itemize}

To jeden z najprostszych sposobow na zachowanie porzadku w projekcie.

\section{O czym warto pamietac na koniec}
Metadata nie sa dodatkiem kosmetycznym. Wplywaja na:

\begin{itemize}
  \item wyglad strony w przegladarce,
  \item SEO,
  \item udostepnianie w mediach spolecznosciowych,
  \item odbior strony na urzadzeniach mobilnych,
  \item dolaczanie zasobow takich jak CSS, favicon czy skrypty.
\end{itemize}

W dobrze zaprojektowanym projekcie MoonChunk metadata sa traktowane jako element
pierwszej klasy, a nie jako cos dopisywanego na koncu.
